\subsection{Checklist: Limitations}
\label{app:checklist_limitations}

\paragraph{Question.}
Does the paper discuss the limitations of the work performed by the authors?

\paragraph{Answer.}
Yes.

\paragraph{Justification.}
The Discussion explicitly identifies what the experiments do and do not establish.
The paper claims that sparse Koopman autoencoders can expose support families that
act as label-free, inspectable regime variables, and it supports this claim with
forecasting, support-coordinate interventions, and support--basin alignment
diagnostics.  It also states that the evidence is strongest in deliberately
controlled settings and should not be read as a proof of a globally smooth
finite-dimensional Koopman linearization for a multibasin state space.

\paragraph{Concrete limitation details and interpretation.}
\begin{itemize}
  \item \textbf{Controlled multibasin scope.}
  The strongest basin-identification evidence comes from procedurally generated
  two-dimensional systems whose attractor geometry is known.  These systems are
  appropriate for isolating whether learned sparse supports can align with
  basin-scale structure, but they do not by themselves establish the same
  behavior for high-dimensional physical systems, partially observed systems, or
  systems whose attractors and basin geometry are unknown.

  \item \textbf{Basin-interior evaluation slice.}
  The main support--basin alignment diagnostic is evaluated on held-out states in
  the per-basin top quartile of a basin-depth margin.  This is an intentionally
  clean slice for asking whether support families identify basin interiors.  It
  does not settle behavior near separatrices, during rare transitions, or in
  regions where small perturbations can change basin assignment or sparse
  support.

  \item \textbf{Evaluation labels are unavailable at deployment.}
  Basin labels, basin counts, and attractor locations are not used to train the
  sparse encoder, fit the Koopman transition, or construct support families in
  the intended setting.  They are used only for benchmark evaluation and
  diagnostic slicing.  Consequently, the reported entropy and family-count
  metrics should be interpreted as post-hoc evidence that the learned supports
  align with known benchmark basins, not as a supervised routing procedure.

  \item \textbf{External chaotic-flow evidence is forecasting-only.}
  The Dysts \(dt{\times}30\) experiments broaden the forecasting evidence beyond
  the procedural generator and show that sparse-latent models remain competitive
  on nontrivial chaotic systems.  They are not basin-identification tests,
  because those systems do not provide the same evaluated basin-label structure
  used for the controlled support--basin diagnostics.

  \item \textbf{Support-family construction has fixed design choices.}
  The main support-family object \(F_{\rm abs}\) starts from an absolute latent
  activation threshold and then merges exact supports using a fixed Jaccard
  threshold.  These choices affect fragmentation and merging: a loose merge can
  collapse distinct regimes, while a strict merge can create many families and
  lower conditional entropy by over-fragmenting basins.  The family count must
  therefore be interpreted together with \(H(B\mid F_{\rm abs})\), rather than
  optimized or read in isolation.

  \item \textbf{Forecasting protocol is a benchmark diagnostic.}
  Forecast horizons and re-encoding periods are measured in stored observation
  steps, and equal horizon values need not correspond to equal physical time
  across systems.  Periodic re-encoding uses only the model's predicted state,
  not the true future or an oracle basin label, but the best-periodic summaries
  should still be read as controlled deployment diagnostics rather than as a
  guarantee of autonomous long-horizon stability in all regimes.
\end{itemize}
